\section{Relation-Level Analysis Framework}
\label{chap:framework}

The previous chapter introduced predictive multiplicity as a way to measure
query-level instability among comparable prediction outputs.
This chapter defines the relation-level analysis framework used in this thesis.
The goal is not only to ask whether predictive multiplicity exists, but also to
identify where it is concentrated in the structure of a knowledge graph.

The basic unit of link prediction evaluation is a query generated from a test
triple.
However, in this thesis, the question is not only whether individual queries are
solved consistently, but also whether such instability is concentrated in
particular relations.
Therefore, we aggregate query-level outcomes into relation-level statistics and
compare these statistics across candidate structural factors.

This chapter first defines the relation-level aggregation procedure, then
introduces the relation-level metrics and the structural factors that are used 
in the experimental chapter.

\subsection{Relation-Level Aggregation}

Let $\mathcal{E}$ be the set of entities and $\mathcal{R}$ be the set of
relations.
A knowledge graph consists of triples $(h,r,t) \in
\mathcal{E}\times\mathcal{R}\times\mathcal{E}$, where $h$ is the head entity,
$r$ is the relation, and $t$ is the tail entity.
In the standard link prediction setting, each test triple $(h,r,t)$ gives
 two prediction queries:
\begin{align}
    q^{\mathrm{head}} = (?,r,t),
    \qquad
    q^{\mathrm{tail}} = (h,r,?).
\end{align}
The first query asks the model to recover the missing head entity, while the
second asks it to recover the missing tail entity.

For each relation $r$, let $\mathcal{Q}^{\mathrm{head}}_r$ and
$\mathcal{Q}^{\mathrm{tail}}_r$ denote the sets of head-prediction and
tail-prediction queries whose relation is $r$, respectively.
We also write
\begin{align}
    \mathcal{Q}_r
    =
    \mathcal{Q}^{\mathrm{head}}_r
    \cup
    \mathcal{Q}^{\mathrm{tail}}_r
\end{align}
for the combined query set of relation $r$.
If relation $r$ appears in $n_r$ test triples, then
$|\mathcal{Q}^{\mathrm{head}}_r|=|\mathcal{Q}^{\mathrm{tail}}_r|=n_r$ and
$|\mathcal{Q}_r|=2n_r$.
We call $n_r$ the test support of relation $r$, because it gives the
number of test triples available for computing relation-level statistics.

The relation-level aggregation used in this thesis computes predictive
performance and multiplicity statistics for each relation separately.
This choice is important for two reasons.
Firstly, it prevents a small number of high-support relations from dominating the
analysis.
Secondly, it makes it possible to compare relations as structural units.
Instead of asking only whether the entire test set has high or low
multiplicity, we ask whether relations with different structural properties
show systematically different multiplicity behavior.

In addition to the combined relation-level view over $\mathcal{Q}_r$, this
thesis keeps the prediction side when it is structurally relevant.
For example, many-to-one and one-to-many mapping patterns have different
implications for head prediction and tail prediction.
Therefore, the main mapping-type analysis in Chapter~\ref{chap:experiments}
uses side-aware relation-level aggregation over
$\mathcal{Q}^{\mathrm{head}}_r$ and $\mathcal{Q}^{\mathrm{tail}}_r$, rather
than relying only on the combined set $\mathcal{Q}_r$.

All ranking results are interpreted under the filtered evaluation convention,
which means that when a query is evaluated, other known true answers for the 
same query are removed before ranking.
Consequently, a lower Hits@10 or higher multiplicity for a structurally harder
prediction direction should not be interpreted simply as an artifact of
multiple true answers competing with each other in an unfiltered candidate set.
On the contrary, it indicates that even after filtering known true answers, 
the ranking decision for that relation direction remains more difficult or less
stable across repeated outputs.

\subsection{Relation-Level Aggregation of Metrics}

Chapter~\ref{chap:predictive_multiplicity} introduces Hits@10, ambiguity, and
discrepancy as metrics for link prediction performance and predictive
multiplicity.
Here, we define how these metrics are aggregated for each relation.

Let $\mathcal{M}$ be the set of repeated prediction outputs used in the
analysis.
In this thesis, these outputs come from repeated runs of the same KGE model
family with different random seeds.
For a query $q$ and an output $m \in \mathcal{M}$, let
\begin{align}
    H_m(q) =
    \begin{cases}
        1, & \text{if output } m \text{ ranks the correct entity within top-}k,\\
        0, & \text{otherwise.}
    \end{cases}
\end{align}
In the experiments, we use $k=10$.

For relation $r$, the relation-level Hits@10 is computed by averaging the
hit indicator over the queries of relation $r$ and over repeated outputs:
\begin{align}
    hits_r
    =
    \frac{1}{|\mathcal{M}|\,|\mathcal{Q}_r|}
    \sum_{m\in\mathcal{M}}
    \sum_{q\in\mathcal{Q}_r}
    H_m(q).
\end{align}
When the prediction side is analyzed separately, $\mathcal{Q}_r$ is replaced by
$\mathcal{Q}^{\mathrm{head}}_r$ or $\mathcal{Q}^{\mathrm{tail}}_r$.

The multiplicity metrics are computed on a relation-specific eligible subset.
For relation $r$, we define
\begin{align}
    \mathcal{Q}^{\mathrm{elig}}_r
    =
    \{q\in\mathcal{Q}_r \mid \max_{m\in\mathcal{M}} H_m(q)=1\}.
\end{align}
This subset contains queries for which at least one repeated output succeeds
within top-$k$.
Queries missed by all outputs are treated as shared failures.
They are not used to measure disagreement among successful and unsuccessful
outputs.

The relation-level ambiguity score is then defined as
\begin{align}
    \alpha_r
    =
    \frac{1}{|\mathcal{Q}^{\mathrm{elig}}_r|}
    \sum_{q\in\mathcal{Q}^{\mathrm{elig}}_r}
    \mathbf{1}
    \left[
        \exists m,m'\in\mathcal{M}
        \text{ such that }
        H_m(q)\neq H_{m'}(q)
    \right].
\end{align}
Thus, $\alpha_r$ measures how often the repeated outputs disagree on the
top-$k$ hit-or-miss decision, among the queries where at least one output is
able to hit.

The relation-level discrepancy score uses the same eligible subset.
Let $m_0$ denote the reference output used in the multiplicity computation.
We define
\begin{align}
    \delta_r
    =
    \max_{m\in\mathcal{M}}
    \frac{1}{|\mathcal{Q}^{\mathrm{elig}}_r|}
    \sum_{q\in\mathcal{Q}^{\mathrm{elig}}_r}
    \mathbf{1}\left[H_m(q)\neq H_{m_0}(q)\right].
\end{align}
This metric measures the largest hit-or-miss disagreement between the reference
output and another repeated output for relation $r$.

The eligible subset is important for interpretation.
Both $\alpha_r$ and $\delta_r$ measure instability among outputs on queries
where success is observed at least once.
They should not be read as error rates over all test queries of relation $r$.
This distinction is used throughout the experimental chapter.

\subsection{Prediction Side}

TODO: Explain why head prediction and tail prediction should be analyzed separately.

\subsection{Candidate Structural Factors}

TODO: Introduce the candidate relation-level structural factors.

\subsubsection{Mapping Type}

TODO: Define one-to-one, one-to-many, many-to-one, and many-to-many mapping types.

\subsubsection{Inverse-Like Support}

TODO: Define the inverse-like metrics used in the thesis, including the stricter high-confidence metrics.

\subsubsection{Symmetry}

TODO: Define symmetry and explain the self-loop caveat.

\subsubsection{Relation Frequency}

TODO: Keep relation frequency as a pending support/sparsity control analysis until its role is confirmed.
